%% kapa: exam=1 round=27 year=2016 period=1 form=A subject=civil_law subject_ko=민법 seq=37 number=37
%% source: Q-Net 감정평가사 기출문제 (cst00309) artlSeq=5169402 file=extracted/1008/5169402/제27회 감정평가사 1차 문제/제27회 감정평가사 1교시문제(A형).pdf
%% conversion: kordoc 4.24.51 -> tools/gen1.py (auto); review=proofread-sonnet; flags=none
\begin{question}{37}
乙은 자기 소유의 돼지 1천 마리를 甲에게 유동집합물로서 점유개정의 방식으로 양도담보한 후 계속 사육하고 있다. 이에 관한 설명으로 옳지 않은 것은? (다툼이 있으면 판례에 따름)
\begin{choices}
  \item 甲의 양도담보권의 효력은 원칙적으로 위 돼지들이 출산한 새끼돼지들에도 미친다.
  \item 만일 화재로 위 돼지들이 폐사하여 乙이 화재보험금청구권을 취득하면, 이에 대해 甲은 양도담보권에 기한 물상대위권을 행사할 수 있다.
  \item 만일 乙이 위 돼지들을 丙에게 점유개정의 방식으로 인도한다면, 丙은 선의취득을 할 수 없다.
  \item 만일 乙이 양도담보 사실을 알고 있는 丁에게 위 돼지들을 양도한 후, 丁이 5백 마리의 돼지를 새로 구입하여 반입한 경우, 별도자금을 투입해 반입한 사실을 증명하면 甲의 양도담보권의 효력은 그 새로 구입한 돼지들에게 미치지 않는다.
  \item 만일 乙이 위 돼지들을 戊에게 점유개정의 방식으로 인도하여 이중으로 양도담보한다면, 戊는 양도담보권을 선의취득한다.
\end{choices}
\end{question}
\answer{5}
